\documentclass{article}
\usepackage{hyperref,amsmath}
\begin{document}
\section*{A receiver-local probe for a changing old-result supplier}
\paragraph{What changes.}
PR 659 required preparation by tick 7 of the one supplier that
could deliver the old keyed result's suffix at tick 8. Both suppliers,
Q and P, held the same exact suffix. Its truth and identity were
therefore insufficient to choose a supplier. The authored message
containing the hidden mode is removed. The preparation actor instead
reads the callback of a scheduled one-byte probe to the very same
\texttt{owed} endpoint, bound to Base in the native replay.
The local binding of callback and actor is an explicit interface:
it does not establish a free remote notification or grant a license.

\paragraph{What is qualified.}
A declared closed probe round finishes before commitment, occupies a
known receiver slot, and excludes unreported dispatch loss, nondelivery
other than the modeled route, and a further route change before the
suffix. Within that family, delivered means the probed route is
usable and a checked absence means the other route is usable.
The candidate still needs an authorized sender, metered time and bytes,
a successful separate preparation command, and the exact matching
prefix and suffix. Loss, late delivery, and competing receiver charge
are tested in envelopes that withhold the closed-round contract;
the policy must treat a missing callback there as unknown, not as
evidence for the alternate. Already receiving a correct short probe
is not generally proof of a future large transfer. This harness
constructs route reachability and the closed-round premise; the
broker does not infer either as a stock sensor guarantee.

\paragraph{Matched resource result.}
In the closed round, prior receiver debit at the probe is
$130+114=244$ bytes. Adding one fits at cap 320; the older 130-byte
item expires before the suffix. With a Q probe, its source debit
even for the suffix route is $140+4+1+173=318$, so the second
four-byte diagnostic still uses its later slot. When P is selected,
Q's unsuccessful probe has still consumed a byte, but both diagnostics
fit. P has no probe debit and its old bound remains
$d+114+173\le320$, or $d\le33$. Probing P instead admits its suffix
only for $d\le32$, while still giving a correct observation under
the closed law. One in-band probe is enough for the binary decision;
zero bits cannot distinguish opposed preparation actions on
identical old receipts at the stated information cut. This is
not a lower bound over changed authority, preprepared redundant
senders, or other independent sensors.

A 76-byte additional arrival at tick 3 makes the receive window
exactly full at tick 7 without consuming otherwise available
suffix capacity at tick 8. Even when the probed sender's path is
clear, its one-byte attempt then receives no callback. Keeping
that case admits an indistinguishable opposite mode, so absence
alone cannot justify binding the alternate. A probe after the cut
cannot repair the missed preparation. An erased preparation request
is also distinct from a correct observation: neither a probe
callback nor a computed selection substitutes for executing the
preparation.

\paragraph{Comparison and scope.}
Replays run pinned SubT dispatch and its observed charge, the existing
RWARE and SQLite handoff fixture, and the unchanged shared admission
checker on the attained observation cells. Unrelated warehouse work
and both diagnostics still finish in the asserted test cases.
A plain direct calculation given the same lawful callback, preparation
power and actual ledger ties. The stable old-result lookup reduces
to scoped retry rendezvous for the digital slice (RIFL, PRs 853--854).
The result is an inexpensive realized observation in a declared
held-in simulator family, not an advance over equally informed
ordinary probing, an independent mission, or a field reliability
claim.
\end{document}
